\documentclass{article}
\usepackage[utf8]{inputenc}
\usepackage{authblk}
\usepackage{graphicx} % Lib de importar e tratar imagem
\usepackage[english,brazil]{babel} % Manter dois idiomas
\usepackage{fancyhdr} % Configuração das paginas
\usepackage{float}
\usepackage{xcolor}
\usepackage{titlesec}
\usepackage{background} % Componente de Watermark
\usepackage{hyperref}
\usepackage{verbatim}
\usepackage[a4paper, left=2cm, right=2cm]{geometry}
\usepackage[most]{tcolorbox}
\usepackage{upquote}
\usepackage{amsmath} % Pacote para fórmulas matemáticas avançadas

\newcommand{\shellcmd}[1]{\\\indent\indent\texttt{\footnotesize\# #1}\\}

% Definir as características do hyperlink
\hypersetup{
    colorlinks=true,
    linkcolor=blue,
    filecolor=magenta,      
    urlcolor=cyan,
    pdftitle={Overleaf Example},
    pdfpagemode=FullScreen,
}

%Função para diminuir tamanho da fonte
\newcommand{\changefont}{%
    \fontsize{9}{12}\selectfont
}

%Document Layout Config
\pagestyle{fancy}
\setlength\headheight{26pt}
\setlength\parindent{0pt}
\fancyhf{}

% Header Config
\fancyhead[RO]{Gerador de Senha}
\fancyhead[LO]{\includegraphics[width=0.7cm]{Logo.png}}

% Footer Config
\renewcommand{\footrulewidth}{0.4pt} % Inserir linha no final da pagina "pre-footer"
\fancyfoot[RO]{\includegraphics[width=0.7cm]{Logo.png}}

%Document Specification
\title{Gerador de Senha}
\author{senhasegura}
\affil{Departamento de Segurança e Inovação - SEGi9}
\date{}
\graphicspath{ {img/} }

%Capa
\begin{document}
\NoBgThispage % Ignorar background na capa
\maketitle
\begin{center}
\includegraphics{logo.png}
\end{center}
\begin{center}
\textbf{\LARGE \\PÚBLICO}
\end{center}
\pagenumbering{gobble} % Não contar como página a capa
\pagebreak

% Background a partir da segunda pagina
\backgroundsetup{contents=\includegraphics{logo.png}, scale=1.5, opacity=0.15, angle=0} % marca d'agua

%Resumo Executivo
\section*{\centering{Resumo Executivo}}
\vspace*{\fill}
\large{As senhas possuem níveis variados de usabilidade e segurança. Senhas simples e fáceis de memorizar são mais propensas a comprometimentos, pois estão vulneráveis a diferentes tipos de ataques. Por isso, é fundamental selecionar senhas que atendam a critérios específicos, incorporando características que aumentem sua resistência a tais ameaças.}
\vspace*{\fill}
\fancyfoot[CO]{senhasegura\\ \changefont Qualquer reprodução ou distribuição é estritamente proibida sem a aprovação prévia por escrito da empresa MT4.}
\pagebreak

%Sumario
\tableofcontents
\pagebreak

%Texto
\fancyfoot[CO]{\thepage\\\changefont Qualquer reprodução ou distribuição é estritamente proibida sem a aprovação prévia por escrito da empresa MT4.}
\pagenumbering{arabic}

% Página
\section{Senha}

A autenticação é um mecanismo de segurança utilizado para verificar a identidade de um indivíduo. Um dos principais métodos de autenticação é o uso de senhas, que, em conjunto com o nome de usuário, permite realizar o processo de login, garantindo acesso à conta e aos dados associados.

\subsection{Características}

Um gerador de senhas possibilita ao usuário configurar diversas características, como a quantidade de caracteres, a inclusão de letras maiúsculas, minúsculas, símbolos e números.\\

Uma senha pode conter as seguintes características:

\begin{itemize}
	\item Letras minúsculas (a até z)
	\item Letras maiúsculas (A até Z)
	\item Caracteres especiais - ! @ \# \$ \% \^{} \& *
	\item Números (0 até 9)
	\item Repetição ou não de caracteres
\end{itemize}

\pagebreak

% Página
\section{Estudo}

A segurança de sistemas computacionais está intimamente ligada à robustez das senhas utilizadas. Este estudo foi desenvolvido com o objetivo de avaliar diferentes características de senhas, dividindo a análise em duas partes principais: senhas com propriedades fixas e senhas com propriedades variadas.

\subsection{Fixa}

Nesta etapa, foram avaliadas senhas que seguem um padrão predefinido de propriedades, incluindo:\\

\textbf{Calcular o número de combinações para cada tipo de caractere.}\\
Calculando a quantidade de caracteres possíveis.
\begin{itemize}
	\item 03 Letras maiúsculas : $26^3$ caracteres
	\item 03 Letras minúsculas : $26^3$ caracteres
	\item 03 Números : $10^3$ caracteres
	\item 03 Caracteres especiais : $8^3$ caracteres
	\item Ordem dos caracteres não importa;
	\item Os mesmos caracteres podem ser usados mais de uma vez;
\end{itemize}

\textbf{Determinar o número de arranjos possíveis}\\
Após determinar os grupos de caracteres, precisamos calcular de quantas maneiras podemos organizar os 12 caracteres (3 maiúsculas, 3 minúsculas, 3 números, e 3 caracteres especiais). Esse é um problema de permutações de multiconjuntos, onde os elementos são indistinguíveis dentro de cada grupo.\\

A fórmula geral para permutações de multiconjuntos é:

\[
P = \frac{n!}{k_1! \cdot k_2! \cdot k_3! \cdot k_4!}
\]

Onde:
\begin{itemize}
    \item $n$ é o número total de elementos,
    \item $k_1, k_2, k_3, k_4$ são os tamanhos dos grupos.
\end{itemize}

Neste caso:
\begin{itemize}
    \item $n$ = 12,
    \item $k_1 = k_2 = k_3 = k_4$  = 3.
\end{itemize}

Substituindo:

\[
P = {\frac{12!}{3! \cdot 3! \cdot 3! \cdot 3!}} { = 369.600}
\]\\

Isso significa que existem 369.600 maneiras de organizar os 12 caracteres.\\

\textbf{Calcular o número total de senhas}\\

O número total de senhas possíveis é o produto das combinações para cada grupo de caracteres e o número de arranjos possíveis:\\

\[
\text{Total de Senhas} = P \cdot \big( (26^3) \cdot (26^3) \cdot (10^3) \cdot (8^3) \big)
\]

\[
\text{Total de Senhas} = 369.600 \cdot 17.576 \cdot 17.576 \cdot 1.000 \cdot 512
\]

\[
\text{Total de Senhas} = 58.457.738.654.515.200.000 = 5,84577386545152 \times 10^{19}
\]

Este é o total de senhas possíveis para as especificações dadas.

\subsection{Variada}

Nesta etapa, a senha é gerada de forma aleatória, sem considerar as características definidas para a Senha Fixa mencionadas no item anterior exceto o tamanho de 12 caracteres.\\

\textbf{Calcular o número de combinações para cada tipo de caractere.}\\
Calculando a quantidade de caracteres possíveis.
\begin{itemize}
	\item Letras maiúsculas : $26^3$ caracteres
	\item Letras minúsculas : $26^3$ caracteres
	\item Números : $10^3$ caracteres
	\item Caracteres especiais :  - $8^3$ caracteres
\end{itemize}

Portanto, o número total de caracteres disponíveis para cada posição na senha é a soma dos tamanhos desses conjuntos:\\

\[
N_{\text{total}} = 26 + 26 + 8 + 10 = 70
\]\\

\textbf{Calcular o número total de senhas}\\

Se tivermos \( n \) caracteres possíveis e uma senha de \( L \) dígitos, o número total de combinações possíveis é dado por:

\[
\text{Total de Senhas} = n^L
\]

Onde:
\begin{itemize}
    \item \( n = 70 \) (número total de caracteres disponíveis),
    \item \( L = 12 \) (comprimento da senha).
\end{itemize}

Substituímos os valores de \( n \) e \( L \) na fórmula:

\[
\text{Total de Senhas} = 70^{12}
\]

Calculando:

\[
\text{Total de Senhas} = 13.841.287.201.000.000.000.000 = 1,3841287201 \times 10^{22}
\]

Este é o total de senhas possíveis para as especificações dadas.

\pagebreak

% Página
\section{Conclusão}

Este estudo apresentou uma análise comparativa sobre o tamanho e as características de senhas, dividido em duas abordagens principais: senhas com propriedades fixas e senhas com propriedades variadas. A primeira abordagem envolveu senhas geradas a partir de critérios predefinidos, resultando em um total de combinações igual a 
\[
\text{Total de Senhas} = 58.457.738.654.515.200.000 = 5,84577386545152 \times 10^{19}
\]

A segunda abordagem considerou uma configuração mais ampla, com variações diversificadas, culminando em 
\[
\text{Total de Senhas} = 13.841.287.201.000.000.000.000 = 1,3841287201 \times 10^{22}
\]
combinações possíveis.\\

Os resultados indicam que senhas com propriedades variadas oferecem um espaço de chave significativamente maior, contribuindo para uma robustez superior contra ataques de força bruta.

\pagebreak

%Copyright
\vspace*{\fill}
\begin{flushright}
	\underline{\textit{Copyright}}\\
	All rights reserved to MT4 Tecnologia LTDA.\bigskip

	\underline{\textit{Propósito deste documento}}\\
	Gerador de Senha\bigskip

	\underline{\textit{Versão}}\\
	Versão 1.0\bigskip

	\underline{\textit{Autor}}\\
	Caio Abreu Ferreira [\href{mailto:cferreira@senhasegura.com}{cferreiras@senhasegura.com}] \bigskip\\
	
	\underline{\textit{MT4 Tecnologia Ltda.}}\\
	Rua Joaquim Antunes, 767 – Pinheiros – São Paulo\\
	Phone: +55 11 3065-7070\\
	\url{www.mt4.com.br}\\
	\href{mailto:support@senhasegura.com}{support@senhasegura.com}\\
\end{flushright}

\end{document}